Queue-worker systems execute heterogeneous operational work such as replication,
repair, cache rebuilds, imports, external-service calls, validation, and media
processing. Profiling such work is difficult because the behavior of a task
attempt is not determined by task kind alone: intrinsic task features and
execution context can change active resource demand, bounded worker occupancy,
tail behavior, and the value of additional instrumentation. Existing metrics,
traces, profiles, and resource-management models provide useful observations,
but they do not by themselves define an attempt-attributed model that separates
task identity, intrinsic work, execution context, cost, uncertainty, risk, and
profiling policy.

This paper develops Task Behavior Graphs, a domain-driven model for profiling
queue-worker task attempts under uncertain resource demand. The model separates
task kind, task instance, and task attempt; represents cost through typed
resource-demand and worker-occupancy dimensions; treats raw observations as
evidence that must be assigned semantic roles; and connects canonical features
to cost estimates, reliability diagnostics, risk-relevant consequences, and
profiling-policy triggers. The contribution is a formal modeling framework and
worked structure for future implementations and evaluations, not an empirical
claim about a particular production system.
